% ===========================================================================
%  Annexe — Glossaire des notions
% ===========================================================================
\chapter{Glossaire des notions}
\markboth{Glossaire}{Glossaire}

Un mot mal compris coûte plus cher qu'une formule oubliée. La formule, on la
retrouve dans le formulaire ; le mot, lui, commande toute la lecture de
l'énoncé. « Bijection », « barycentre », « équiprobable », « congru » : ces
termes ont, en mathématiques, un sens étroit, et c'est ce sens-là que le
correcteur attend.

Les pages qui suivent rassemblent les termes de l'année dans l'ordre
alphabétique. Chaque définition est courte à dessein : elle rappelle ce que
le mot désigne, non comment on s'en sert. Le numéro du chapitre indiqué à
droite renvoie à l'endroit où la notion est construite, démontrée et mise au
travail.

% ---------------------------------------------------------------------------
% Mise en forme locale
\newcommand{\glolettre}[1]{%
  \par\vspace{1.15em}\noindent
  \begingroup\sffamily\bfseries\large\color{bmtAccent}#1\endgroup
  \hspace{0.5em}%
  {\color{bmtGrisClair}\leaders\hrule height 1pt\hfill\kern0pt}%
  \par\vspace{0.35em}}

\newcommand{\glo}[2]{%
  \par\vspace{0.5em}\noindent
  {\sffamily\bfseries\small\color{bmtAccent}#1}%
  \hspace{0.45em}{\sffamily\footnotesize\color{bmtGris}ch.~#2}%
  \par\vspace{0.05em}\noindent\ignorespaces}

% ---------------------------------------------------------------------------
\glolettre{A}

\glo{Affixe}{12}
Nombre complexe associé à un point ou à un vecteur du plan : le point
$M(a\,;b)$ a pour affixe $z = a+ib$. C'est le pont qui fait passer d'une
figure à un calcul.

\glo{Antidéplacement}{14}
Isométrie du plan qui renverse l'orientation : les symétries axiales et les
symétries glissées, et elles seules.

\glo{Argument}{12}
Mesure de l'angle que fait avec l'axe des abscisses le vecteur d'affixe $z$
non nulle. Il est défini à $2k\pi$ près, ce qui oblige à préciser chaque fois
l'intervalle choisi.

\glo{Arrangement}{17}
Liste ordonnée de $p$ éléments distincts pris parmi $n$. L'ordre compte : deux
arrangements formés des mêmes éléments rangés autrement sont différents.

\glo{Asymptote}{3}
Droite dont la courbe se rapproche indéfiniment. Verticale quand la fonction
explose en un point, horizontale ou oblique quand la différence entre la
courbe et la droite tend vers zéro à l'infini.

\glo{Automorphisme}{11}
Application linéaire d'un ensemble dans lui-même qui est bijective. En
matrices, cela revient à un déterminant non nul.

% ---------------------------------------------------------------------------
\glolettre{B}

\glo{Barycentre}{13}
Point d'équilibre d'un système de points affectés de coefficients. Il n'existe
que si la somme des coefficients n'est pas nulle, et il ne change pas quand on
multiplie tous les coefficients par un même réel non nul.

\glo{Bijection}{1, 5}
Fonction qui, à chaque valeur de l'ensemble d'arrivée, fait correspondre un
antécédent et un seul. Une fonction continue et strictement monotone sur un
intervalle en réalise toujours une.

\glo{Binomiale (loi)}{18}
Loi de probabilité du nombre de succès obtenus en répétant $n$ fois, de façon
indépendante, une même épreuve à deux issues.

% ---------------------------------------------------------------------------
\glolettre{C}

\glo{Combinaison}{17}
Sous-ensemble de $p$ éléments pris parmi $n$, sans considération d'ordre. Le
nombre de ces sous-ensembles est noté $\combi{p}{n}$.

\glo{Congruence}{9}
Relation $a \equiv b \pmod n$ signifiant que $a-b$ est un multiple de $n$.
Elle permet de raisonner sur les restes plutôt que sur les nombres.

\glo{Conjugué}{12}
Nombre complexe $\overline z = a-ib$ associé à $z = a+ib$. Géométriquement,
c'est le symétrique de $z$ par rapport à l'axe des réels.

\glo{Continuité}{1}
Propriété d'une fonction dont la limite en un point est égale à sa valeur en
ce point. Sur un intervalle, elle traduit une courbe que l'on trace sans lever
le crayon.

\glo{Convexité}{2}
Une fonction est convexe sur un intervalle si sa courbe y reste au-dessus de
chacune de ses tangentes ; cela équivaut à une dérivée seconde positive.

\glo{Coplanaires}{16}
Se dit de points ou de vecteurs appartenant à un même plan. Trois vecteurs
sont coplanaires exactement lorsque leur déterminant est nul.

% ---------------------------------------------------------------------------
\glolettre{D}

\glo{Déplacement}{14}
Isométrie du plan qui conserve l'orientation : translations et rotations, et
elles seules.

\glo{Dérivabilité}{2}
Existence d'une limite finie du taux d'accroissement en un point. Cette limite
est le nombre dérivé, et c'est le coefficient directeur de la tangente.

\glo{Déterminant}{11, 16}
Nombre attaché à une matrice carrée, ou à un triplet de vecteurs de l'espace.
Sa nullité signale une dépendance : matrice non inversible, vecteurs
coplanaires.

\glo{Diviseur}{9}
Un entier $d$ divise $a$ s'il existe un entier $k$ tel que $a = kd$. La
division euclidienne, elle, produit toujours un quotient et un reste, que la
division soit exacte ou non.

% ---------------------------------------------------------------------------
\glolettre{E}

\glo{Écart type}{18}
Racine carrée de la variance. Il mesure la dispersion des valeurs autour de
l'espérance, dans la même unité que la variable.

\glo{Équation différentielle}{7}
Équation dont l'inconnue est une fonction et qui relie cette fonction à ses
dérivées. Résoudre, c'est décrire toutes les fonctions qui la vérifient.

\glo{Équiprobabilité}{17}
Situation où toutes les issues d'une expérience ont la même probabilité. La
probabilité d'un événement se réduit alors au quotient du nombre de cas
favorables par le nombre de cas possibles.

\glo{Espérance}{18}
Moyenne des valeurs d'une variable aléatoire, chacune pesée par sa
probabilité. C'est le gain moyen que l'on obtiendrait en répétant très
longtemps l'expérience.

\glo{Événement}{17}
Partie de l'univers, c'est-à-dire ensemble d'issues. Il est réalisé dès que
l'issue observée lui appartient.

\glo{Exponentielle}{6}
Fonction réciproque du logarithme népérien, notée $\exp$ ou $x \mapsto e^{x}$.
Elle est sa propre dérivée, ce qui explique sa présence dans toutes les
équations différentielles du programme.

% ---------------------------------------------------------------------------
\glolettre{F}

\glo{Fonction réciproque}{5, 6}
Fonction qui défait ce que fait une bijection : si $f(a) = b$, alors
$f^{-1}(b) = a$. Les deux courbes sont symétriques par rapport à la droite
d'équation $y = x$.

% ---------------------------------------------------------------------------
\glolettre{H}

\glo{Homothétie}{14, 15}
Transformation qui, à partir d'un centre $\Omega$ et d'un rapport $k$ non nul,
envoie $M$ sur $M'$ tel que $\vect{\Omega M'} = k\,\vect{\Omega M}$. Elle
multiplie les longueurs par $\abs k$ et conserve les angles.

% ---------------------------------------------------------------------------
\glolettre{I}

\glo{Indépendance}{17, 18}
Deux événements sont indépendants lorsque la réalisation de l'un ne change pas
la probabilité de l'autre, ce qui se traduit par
$\proba{A \cap B} = \proba A \times \proba B$.

\glo{Intégrale}{4}
Pour une fonction continue et positive, aire du domaine compris entre la
courbe, l'axe des abscisses et deux droites verticales. Dans le cas général,
c'est la variation d'une primitive entre les deux bornes.

\glo{Isométrie}{14}
Transformation du plan qui conserve les distances. Elle conserve donc aussi
les angles géométriques, les milieux et les aires.

% ---------------------------------------------------------------------------
\glolettre{L}

\glo{Limite}{1}
Valeur dont $f(x)$ se rapproche autant qu'on veut lorsque $x$ s'approche d'un
point ou tend vers l'infini. Elle peut être finie, infinie, ou ne pas exister.

\glo{Logarithme népérien}{5}
Unique primitive de $x \mapsto \frac1x$ sur $\Rps$ qui s'annule en $1$. Il
transforme les produits en sommes, ce qui est la raison de son invention.

% ---------------------------------------------------------------------------
\glolettre{M}

\glo{Matrice}{11}
Tableau rectangulaire de nombres. Elle sert ici à écrire d'un trait un système
linéaire, une transformation du plan, ou une suite de relations de récurrence.

\glo{Module}{12}
Distance de l'origine au point d'affixe $z$ : $\abs z = \sqrt{a^{2}+b^{2}}$
pour $z = a+ib$. Le module d'un produit est le produit des modules.

\glo{Monotone}{2, 3}
Se dit d'une fonction croissante sur tout un intervalle, ou décroissante sur
tout un intervalle. La stricte monotonie interdit les paliers.

% ---------------------------------------------------------------------------
\glolettre{N}

\glo{Nombre premier}{9}
Entier supérieur à $1$ n'ayant pour diviseurs positifs que $1$ et lui-même.
Tout entier supérieur à $1$ se décompose en produit de nombres premiers, et
d'une seule façon.

\glo{Numération (base)}{10}
Écriture d'un entier comme somme de puissances d'un nombre $b \geq 2$ appelé
base. Les chiffres utilisés sont alors les entiers de $0$ à $b-1$.

% ---------------------------------------------------------------------------
\glolettre{P}

\glo{PGCD}{9}
Plus grand commun diviseur de deux entiers. L'algorithme d'Euclide le fournit
par divisions successives, sans qu'on ait besoin de décomposer les nombres.

\glo{PPCM}{9}
Plus petit commun multiple de deux entiers. Pour deux entiers positifs $a$ et
$b$, on a toujours $\pgcd(a,b) \times \ppcm(a,b) = ab$.

\glo{Primitive}{4}
Fonction dont la dérivée est la fonction de départ. Deux primitives d'une même
fonction sur un intervalle ne diffèrent que d'une constante.

\glo{Probabilité conditionnelle}{17}
Probabilité d'un événement $B$ sachant qu'un événement $A$ de probabilité non
nulle est réalisé, notée $\probacond BA$. Elle se lit directement sur une
branche d'arbre pondéré.

% ---------------------------------------------------------------------------
\glolettre{R}

\glo{Récurrence}{8}
Mode de démonstration en deux temps : vérifier la propriété au premier rang,
puis montrer qu'elle se transmet d'un rang au suivant.

\glo{Rotation}{14}
Isométrie déterminée par un centre et un angle. Elle conserve les distances et
l'orientation, et ne laisse fixe que son centre — sauf si l'angle est nul.

% ---------------------------------------------------------------------------
\glolettre{S}

\glo{Similitude}{15}
Transformation qui multiplie toutes les distances par un même rapport $k>0$.
Directe si elle conserve l'orientation, indirecte sinon ; les isométries sont
les similitudes de rapport $1$.

\glo{Suite adjacente}{8}
Deux suites, l'une croissante, l'autre décroissante, dont la différence tend
vers zéro. Elles convergent alors vers une même limite, encadrée par leurs
termes.

\glo{Suite arithmétique}{8}
Suite dont on passe d'un terme au suivant en ajoutant toujours le même nombre,
appelé raison.

\glo{Suite géométrique}{8}
Suite dont on passe d'un terme au suivant en multipliant toujours par le même
nombre non nul, appelé raison.

\glo{Symétrie glissée}{14}
Composée d'une symétrie axiale et d'une translation de vecteur dirigeant
l'axe. C'est l'antidéplacement sans point fixe.

% ---------------------------------------------------------------------------
\glolettre{T}

\glo{Tangente}{2}
Droite qui approche la courbe au plus près au voisinage d'un point. Son
coefficient directeur est le nombre dérivé en ce point.

\glo{Théorème des valeurs intermédiaires}{1}
Une fonction continue sur un intervalle prend toute valeur comprise entre deux
de ses valeurs. Avec la stricte monotonie, la solution obtenue est unique.

\glo{Translation}{14}
Isométrie qui déplace tous les points d'un même vecteur. Elle n'a aucun point
fixe, sauf si le vecteur est nul.

% ---------------------------------------------------------------------------
\glolettre{U}

\glo{Univers}{17}
Ensemble de toutes les issues possibles d'une expérience aléatoire. Le choisir
correctement est la première décision d'un exercice de probabilités, et
souvent la seule difficile.

% ---------------------------------------------------------------------------
\glolettre{V}

\glo{Variable aléatoire}{18}
Fonction qui associe un nombre à chaque issue d'une expérience aléatoire. Sa
loi donne, pour chaque valeur prise, la probabilité correspondante.

\glo{Variance}{18}
Moyenne des carrés des écarts à l'espérance. Elle mesure la dispersion, mais
dans une unité au carré, d'où le passage à l'écart type.

\glo{Vecteur normal}{16}
Vecteur non nul orthogonal à tous les vecteurs d'un plan. Ses coordonnées se
lisent directement sur une équation cartésienne du plan.

\glo{Vecteur directeur}{16}
Vecteur non nul dont la direction est celle d'une droite. Deux droites sont
parallèles exactement lorsque leurs vecteurs directeurs sont colinéaires.
